\section{Discussion}\label{sec:discussion}

\subsection{What the results say}

\paragraph{Computing outside the agent model was associated with large gains where the agent model was weak.} The
accuracy gains in \cref{tab:accuracy} concentrate in models that were weak with raw context, and the per-question
results are consistent with an explanation: the questions those models got wrong in \CA{} and \CB{} need arithmetic
over many days, a fitted trend, or knowledge that two sources report the same steps, and a packet hands the agent
model a computed value with its caveats. The study does not isolate this cause. \CC{} bundles computed values,
data-quality caveats, a second, larger model as planner, and a tool description that tells the agent model to rely
on the numbers. The decomposition in \cref{sec:res:failures} shows that runs with a released packet were nearly all correct and that almost no run was left without a packet, which fits the explanation but is not an
ablation. The pattern agrees with
the program-aided reasoning literature~\citep{gao2022pal,chen2023pot} and with PHIA's finding that computing in code
beats reasoning over the numbers~\citep{merrill2024phia}; in \sys{} the agent model writes no code at all.

\paragraph{For models near ceiling, the benefit is not accuracy.} The strongest models answered almost every
question correctly in every condition, and we claim no accuracy gain for them; the differences are at most a few points. What \sys{} changes for these models is what they are given and what remains afterwards. The agent model
read \ctxRatioCAMin{}--\ctxRatioCAMax{} of its raw-context input; the release was limited to the level the question
needed, so a scheduling assistant receives a bucket rather than a night's sleep in minutes; and every release and
denial left a hash-chained receipt. For \code{claude-fable-5.2} and \code{gemini-4-neon} it also lowered the cost per
correct answer; for \code{gpt-6.2-astra} it raised it, and it added latency for every fast model.

\paragraph{Exposure and leakage are different quantities.} Exposure is fixed by the system before the agent model
runs, so it can be measured, audited and bounded; leakage depends on the model, the prompt and the adversary, and
prompt-level instructions reduce it without bounding it~\citep{li2026toolminimize,dong2026claw}. Our privacy results
show both sides. Where the request asked for data outside the role, raw context and open tools leaked often and
\sys{} did not, because the data was never delivered. Where the user asked for their own data to be used, as in the
manager message, the difference came from the release level: the \CC{} agent model received a bucket and had no
figure to restate. Upstream minimization has worked in other settings~\citep{bagdasarian2024airgap,abdelnabi2025firewalls,li2026toolminimize};
\sys{} applies it with computed release levels rather than a second model. Minimal answers can still compose into a
reconstruction of the record~\citep{dinur2003revealing,patil2025sumleaks}; the gate can cap releases per day, but
this study did not exercise that limit or run a reconstruction attack.

\paragraph{Measuring leakage needs a validated detector.} The first leak rule we wrote had a precision of
\auditPrec{} against hand labels in the development run: it read advice to the user, a sync age or a coincidental
number as a leak. A study that reported it alone would have misstated the difference between the designs. We
recommend that agent-privacy evaluations report detector precision and recall on a labelled sample, and freeze the
detector before the runs it scores, as we did.

\paragraph{Where errors remain.} On the synthetic questions, \sys{}'s remaining errors were almost all the agent
model's reading of a correct packet: a partial total reported despite a caveat that it is incomplete, or \emph{yes}
to a threshold when every value received was above it. Planner faults that dominated the development run, a schema
that rejected ``hours'' and a keyword check that read ``period'' as cycle tracking, were found through the reasons the
broker logs for each refusal and did not recur after they were fixed.

\paragraph{What PHIA shows about the design.} PHIA exposed two decisions that cost accuracy on questions the broker
could answer. First, the gate releases the lowest level that can answer an \emph{operation}; for a threshold check
that is a yes-or-no flag, so a question asking \emph{how many} days met a condition received a flag although the
count was computed. The release level should follow the question asked, not only the operation. Second, a metric
outside the grant is denied without saying which granted metric could answer; the planner mapped ``heart rate'' to
the heart-rate metric, the grant held resting heart rate, and the agent model was left with a denial. A denial that
names an answerable alternative keeps the gate's guarantee and lets the agent model recover; in the second
iteration a quarter of the denied runs did. Both changes are
general, not specific to PHIA. We made them after the study and measured them on fresh questions
(\cref{sec:res:v2}): counts now reach the agent model, a quarter of the denied runs recovered through the hint, and
pooled accuracy on answerable questions, as worded, moved level with raw context, though the gap on PHIA's labels grew, while the strongest open models still did better
with raw context and a prompt change cost accuracy on one synthetic question; a third revision that gave the
planner the comparison windows removed that loss and, in an analysis we did not pre-register, raised accuracy on real data (\cref{sec:res:v3}). Two further observations are properties of the benchmark: some reference labels count time in bed or total
sleep where the question says sleep or light sleep, and ``last $N$ days'' is labelled as $N{+}1$ days. We report
PHIA as worded and with its labels for that reason. PHIA also shows the trade-off we chose: \sys{} refused
unsupported questions far more often than the baselines, and refused some answerable ones.

\paragraph{The open-tools condition depends on how the model calls tools.} Open tools were the worst condition for
some models and the best for others. One model wrote most of its tool calls as text, which a harness that only
reads structured calls would have scored as unanswered; the harness executes such calls in every condition, and
with it the model's open-tools accuracy was \accLlamaScoutB{}. A comparison of MCP-style designs therefore depends on the serving stack as much as on the design,
and studies should state how text-format tool calls are handled.

\subsection{Implications for agent builders}

For people building agents over personal data, we draw four lessons. First, give agent models computed answers
with caveats rather than records; the packet contract cost the strongest models nothing measurable in accuracy and
helped the weaker ones. Second, decide disclosure in code, outside the model, and bind purpose at consent time rather
than in the request. Third, log both releases and denials with their reasons; the logged reasons let us find and fix
the planner defects. Fourth, choose the release level from the question and make denials point to what the grant
does allow; a minimal answer that does not answer the question is not minimal, only unhelpful.

\subsection{Limitations of the design}

The gate cannot stop an agent model from misreading a packet, and a released packet is outside the system's
control once an agent's provider stores it. Within the ceiling the grant sets, what is requested is chosen by the
planner and the agent model, so the release is deterministic in its limits, not in its content. The operation set is
deliberately narrow, so \sys{} refuses questions a code-writing agent could answer (\phiaUnsupported{} of PHIA's
\phiaTotal{} objective questions). The threat model assumes the hosting platform is not compromised.
